% This document is compiled using pdfLaTeX
% You can switch XeLaTeX/pdfLaTeX/LaTeX/LuaLaTeX in Settings

\documentclass{article}
\usepackage{ctex}
\usepackage{amsmath} % 确保引入了此宏包以支持 aligned

\title{因子日历2026}

\date{\today}

\begin{document}

\maketitle

\section{LSV 模型(行为金融因子-羊群效应)}

\subsection{因子定义}
\begin{equation}
    \mathrm{LSV}_{i,t} = |p_{i,t} - p_t| - E|p_{i,t} - p_t|
\end{equation}
\begin{equation}
    E|p_{i,t} - p_t| = \mathrm{AF}(i,t) = \sum_{k=1}^{N_{i,t}} \binom{N_{i,t}}{k} p_t^k (1 - p_t)^{N_{i,t} - k} \left| \frac{k}{N_{i,t}} - p_t \right|
\end{equation}
\begin{equation}
    p_{i,t} = \frac{B(i,t)}{B(i,t) + S(i,t)}, \quad p_t = \frac{\sum_{i=1}^N p_{i,t}}{N}, \quad N_{i,t} = B(i,t) + S(i,t)
\end{equation}

\subsection{变量定义}
\begin{itemize}
    \item ① \( B(i,t), S(i,t) \) 分别为区间 \( t \) 内买入股票和卖出股票的总数量；
    \item ② \( N_{i,t} \) 为持仓发生变动的基金数量；
    \item ③ \( p_{i,t} \) 为股票 \( i \) 在区间 \( t \) 内买入股票基金占持仓变动基金的比例；
    \item ④ 此处使用的是基金年报和半年报持仓数据，假设有 300 只基金加仓股票 \( i \)，则 \( B(i,t) = 300 \)；
    \item ⑤ \( N \) 为全市场股票总数。
\end{itemize}

\subsection{说明}
LSV 模型从基金持股变动数据出发判断一只股票是否被集中买入或集中卖出，捕捉基金“抱团”现象。基金“抱团”现象是股票出现羊群效应的潜在驱动因子之一。

\subsection{参考文献}
\begin{enumerate}
    \item Lakonishok, J., Shleifer, A., \& Vishny, R. W. (1992). \textit{The Impact of Institutional Trading on Stock Prices}. Journal of Financial Economics, 32(1), 23-43.
\end{enumerate}

\section{尾盘成交额占比（高频因子-成交分布类）}

\subsection{因子定义}
\begin{equation}
    \mathrm{APL}_{t,d} = \frac{\sum_{i=240-t}^{240} \mathrm{Amount}_{i,d}}{\mathrm{TotalAmount}_d}
\end{equation}
\begin{equation}
    \mathrm{APL}_t = \frac{\sum_{d=1}^{N} \mathrm{APL}_{t,d}}{N}
\end{equation}

\subsection{变量定义}
\begin{itemize}
    \item ① \( \mathrm{Amount}_{i,d} \) 为交易日 \( d \) 内第 \( i \) 分钟的成交额；
    \item ② \( t \) 为尾盘最后的 \( t \) 分钟，如 \( t=30 \) 对应尾盘 30 分钟；
    \item ③ \( N \) 为日频以上的时间区间内的交易日天数，如 \( N=20 \) 对应的月度因子。
\end{itemize}

\subsection{说明}
尾盘成交额占比比早盘成交额占比有更强的收益预测能力。与未来收益负相关，尾盘成交额占比高的股票存在知情交易者比例更多，资金获得利好消息后选择在噪声交易者更少的尾盘买入；取值高可识别为交易驱动型股票，反之为基本面驱动型股票。

\subsection{参考文献}
\begin{enumerate}
    \item 尹治柱. 2020. 高频视角下成交额蕴藏的 Alpha. 华安证券.
\end{enumerate}

\section{多层次订单失衡(高频因子-流动性类)}

\subsection{因子定义}

\begin{equation}
    \mathrm{MOFI}_{\mathrm{weight}} = \frac{\sum_{i=1}^{5} w_i \times \mathrm{OFI}_t^{(i)}}{\sum_{i=1}^{5} w_i}, \quad w_i = \frac{i}{5}, \quad i = 1, 2, 3, 4, 5
\end{equation}

\begin{equation}
    \Delta V_t^B = 
    \begin{cases} 
        -V_{t-1}^B, & P_t^B < P_{t-1}^B \\ 
        V_t^B - V_{t-1}^B, & P_t^B = P_{t-1}^B \\ 
        V_t^B, & P_t^B > P_{t-1}^B 
    \end{cases}
\end{equation}

\begin{equation}
    \Delta V_t^A = 
    \begin{cases} 
        V_t^A, & P_t^A < P_{t-1}^A \\ 
        V_t^A - V_{t-1}^A, & P_t^A = P_{t-1}^A \\ 
        -V_{t-1}^A, & P_t^A > P_{t-1}^A 
    \end{cases}
\end{equation}

\begin{equation}
    \mathrm{OFI}_t = \Delta V_t^B - \Delta V_t^A
\end{equation}

\subsection{变量定义}
\begin{itemize}
    \item ① \( \mathrm{OFI}_t^{(i)} \) 为第 \( i \) 档下的订单失衡因子；
    \item ② \( P_t^B \) 和 \( P_t^A \) 分别为 \( t \) 时刻的买一价和卖一价；
    \item ③ \( V_t^B \) 和 \( V_t^A \) 分别为 \( t \) 时刻的买一量和卖一量。
\end{itemize}

\subsection{说明}
MOFI 衡量了不同档位订单失衡的加权累积影响，能更准确量化订单失衡对股价的短期和长期影响。短期内，订单失衡通常与未来收益正相关；中长期内，随着买卖压力失衡的消失，股票的超额收益出现均值回复，与收益相关性会减弱。

\subsection{参考文献}
\begin{enumerate}
    \item 丁鲁明, 陈升锐. 2021. 多层次订单失衡及订单斜率因子. 中信建投证券.
\end{enumerate}

\section{趋势资金相对均价因子(量价因子改进)}

\subsection{因子定义}

1. 每个交易日 \( t \)，回看过去 \( k \) 个交易日 \( t-k, t-(k-1), \ldots, t-1 \)，计算过去 \( k \) 日分钟成交量的 \( m\% \) 分位数，定义为成交量阈值，\( k=5, m=90 \)；

2. 将交易日 \( t \) 的每分钟成交量与上述阈值进行对比，若某一分钟的成交量大于上述阈值，则该分钟的交易被定义为趋势资金的交易；

3. 每个交易日，计算如下因子：
\begin{equation}
    \frac{\text{趋势资金的成交量加权平均价格}}{\text{所有分钟的成交量加权平均价格}} - 1
\end{equation}

4. 每月月底回看过去 20 个交易日，计算上述因子的平均值，并做横截面市值、中信一级行业中性化处理。

\subsection{说明}
趋势资金相对均价因子值较大，则表明趋势资金的交易价格相对较高，其出货的可能性更大，后市看空；若因子值较小，则表明趋势资金的交易价格较低，其交易行为更有可能是逢低买入，后市看涨。

\subsection{参考文献}
\begin{enumerate}
    \item 沈亚琦, 刘富兵. 2024. 基于趋势资金日内交易行为的选股因子. 国盛证券.
\end{enumerate}

\section{修正的净流入率(高频因子-资金流类)}

\subsection{因子定义}

1. 分别计算主力资金的买入金额 \( B_P \) 和卖出金额 \( S_P \)，即 Wind 的超大单、大单和中单的合并定义为具有市场定价能力的“广义主力资金”：
\begin{equation}
    B_P = B_{\mathrm{Extra}} + B_{\mathrm{Large}} + B_{\mathrm{Med}}
\end{equation}
\begin{equation}
    S_P = S_{\mathrm{Extra}} + S_{\mathrm{Large}} + S_{\mathrm{Med}}
\end{equation}

2. 利用日涨跌幅来修正主力资金的买入金额和卖出金额：
\begin{equation}
    \ln(B_P / S_P) = \alpha + \beta \times \mathrm{Ret} + \varepsilon
\end{equation}
\begin{equation}
    \hat{B}_P = \frac{e^{\varepsilon}}{1 + e^{\varepsilon}} \times (B_P + S_P)
\end{equation}
\begin{equation}
    \hat{S}_P = \frac{1}{1 + e^{\varepsilon}} \times (B_P + S_P)
\end{equation}

3. 计算修正后的主力资金的净流入占比 \( \mathrm{CNIR} \)：
\begin{equation}
    \mathrm{CNIR} = \frac{\sum_{i=1}^N (\hat{B}_{P,i} - \hat{S}_{P,i})}{\sum_{i=1}^N (\hat{B}_{P,i} + \hat{S}_{P,i})}
\end{equation}

\subsection{说明}
从资金流动力学的角度理解，大单资金流与涨跌幅呈正相关缘自主力资金买卖的不平衡，通过选取买入卖出金额比作为代理变量（IMB），并做截面回归的方法，可以剥离反转因素的影响以提纯资金流的 Alpha 信息；具有定价权的主力资金并非仅是 20 万元以上的大单，基于小金额标准（如 2 万元）识别的主力资金因子表现更好。

\subsection{参考文献}
\begin{enumerate}
    \item 魏建榕, 苏良. 2022. 大小单重定标与资金流因子改进. 开源证券.
\end{enumerate}

\section{上行已实现波动率(高频因子-波动跳跃类)}

\subsection{因子定义}
\begin{equation}
    \mathrm{RS}_t^+ = \sum_{i=1}^N r_{t,i}^2 \cdot \mathrm{I}_{r_{t,i} \geq 0}
\end{equation}

\subsection{变量定义}
\begin{itemize}
    \item ① \( r_{t,i} \) 为某股票交易日 \( t \) 内的分钟对数收益率序列，如 5 分钟对数收益率序列；
    \item ② \( \mathrm{I}_{r_{t,i} \geq 0} \) 为示性函数，对数收益率大于等于 0 时取 1，否则取 0；
    \item ③ \( N \) 表示该股票在交易日 \( t \) 中特定频率的分钟级别数据个数，如 5 分钟频率通常有 48 个数据点。
\end{itemize}

\subsection{说明}
上行已实现波动率对应资产价格波动中的上行波动；上行波动率较高的股票短期股价出现大幅拉升，未来更有可能补跌，收益出现反转。

\subsection{参考文献}
\begin{enumerate}
    \item Barndorff-Nielsen, O. E., Kinnebrock, S., \& Shephard, N. (2010). \textit{Measuring downside risk: realised semivariance}. In T. Bollerslev, J. Russell, \& M. Watson (Eds.), Volatility and Time Series Econometrics: Essays in Honor of Robert F. Engle (pp. 117-136). Oxford University Press.
\end{enumerate}

\section{客户接近中心性(图谱网络-网络结构)}

\subsection{因子定义}
\begin{equation}
    c_i = \frac{1}{\frac{1}{N-1} \sum_{j(\neq i)} D_{ij}}
\end{equation}

\subsection{变量定义}
\begin{itemize}
    \item ① \( D_{ij} \) 为供应商公司 \( j \) 到客户公司 \( i \) 的最短路径；
    \item ② \( N \) 为与客户公司 \( i \) 相连的所有供应商公司的总数。
\end{itemize}

\subsection{说明}
接近中心性衡量了供应链网络中能达到某个客户节点的所有供应商节点的平均最短距离的倒数，接近中心性越高，说明该企业与网络中其他企业建立连接的路径长度越短；供应链网络中占据中心位置的企业更能吸引投资者的注意力，市场对其信息反应更快，有较低的股票回报可预测性；中心性较低的企业由于信息传播较慢，有较高的股票回报可预测性。

\subsection{参考文献}
\begin{enumerate}
    \item Bozok, İ., \& Özyıldırım, S. (2022). \textit{Firm centrality and limited attention}. International Review of Economics \& Finance, 78, 483-500.
\end{enumerate}

\section{个人投资者交易比-惊恐因子(高频因子-动量反转类)}

\subsection{因子定义}

1. 计算各股票的日度“惊恐度”：
\begin{equation}
    \text{惊恐度} = \frac{\text{偏离度}}{\text{基准项}} = \frac{\mathrm{abs}(\text{个股收益率} - \text{市场收益率})}{\mathrm{abs}(\text{个股收益率}) + \mathrm{abs}(\text{市场收益率}) + 0.1}
\end{equation}

2. 计算各股票的个人投资者交易比：
\begin{equation}
    \text{个人投资者交易比} = \frac{\text{成交额小于4万元的个人投资者买入和卖出金额均值}}{\text{个股总成交额}}
\end{equation}

3. 对各股票，将每日惊恐度、每日收益率、个人投资者交易比相乘，得到加权调整收益率；

4. 对各股票，每月末计算过去20日加权调整收益率序列的均值和标准差，分别称为“个人投资者交易比-惊恐收益”因子和“个人投资者交易比-惊恐波动”因子，再将二者等权合成为“个人投资者交易比-惊恐”因子。

\subsection{变量定义}
\begin{itemize}
    \item ① 选取中证全指 000985 代表市场水平；
    \item ② “偏离度”衡量了个股收益率对市场收益率的偏离水平；“基准项”衡量了市场总体收益率水平。
\end{itemize}

\subsection{说明}
个人投资者交易比-惊恐因子在原始惊恐因子上新增了个人投资者交易比重，因为“显著效应（投资者会被具有显著收益的股票所吸引，导致它们被错误定价）在小市值股票上更为强烈”，而小市值股票的个人投资者交易占比相对更高，更易受情绪影响而出现反应过度。

\subsection{参考文献}
\begin{enumerate}
    \item 曹春晓, 温晋欣. 极端收益扭曲决策权重和“草木皆兵”因子. 2022, 方正证券.
    \item Cakici, N., \& Zaremba, A. (2022). \textit{Salience theory and the cross-section of stock returns: International and further evidence}. Journal of Financial Economics, 146(2), 689-725.
\end{enumerate}

\section{价量相关性因子(高频因子-量价相关性类)}

\subsection{因子定义}

① 将 \( \mathrm{PV\_corr\_avg} \) 和 \( \mathrm{PV\_corr\_std} \) 分别剔除反转因子（以过去 20 日收益率），再各自横截面标准化，等权线性相加构建综合因子，将得到新因子 \( \mathrm{PV\_corr\_deRet20} \)：
\begin{equation}
    \mathrm{PV\_corr\_deRet20} = \frac{\mathrm{PV\_corr\_avg}_{\perp} - \mathrm{mean}(\mathrm{PV\_corr\_avg}_{\perp})}{\mathrm{std}(\mathrm{PV\_corr\_avg}_{\perp})} 
\end{equation}
\begin{equation}
    +\frac{\mathrm{PV\_corr\_std}_{\perp} - \mathrm{mean}(\mathrm{PV\_corr\_std}_{\perp})}{\mathrm{std}(\mathrm{PV\_corr\_std}_{\perp})}
\end{equation}

② 将 \( \mathrm{PV\_corr\_trend} \) 带来的增量信息叠加到 \( \mathrm{PV\_corr\_deRet20} \) 之上，得到 \( \mathrm{cpv} \)：
\begin{equation}
    \mathrm{cpv} = \frac{\mathrm{PV\_corr\_deRet20} - \mathrm{mean}(\mathrm{PV\_corr\_deRet20})}{\mathrm{std}(\mathrm{PV\_corr\_deRet20})} + \frac{\mathrm{PV\_corr\_trend} - \mathrm{mean}(\mathrm{PV\_corr\_trend})}{\mathrm{std}(\mathrm{PV\_corr\_trend})}
\end{equation}

\subsection{说明}
价量相关性因子综合了价量相关性平均强度、波动性、趋势性 3 方面的信息，对传统反转因子进行了修正，具备良好的选股能力。

\subsection{参考文献}
\begin{enumerate}
    \item 高子剑. 2020. 高频价量相关性，意想不到的选股因子. 东吴证券.
\end{enumerate}

\section{公司特征相似度动量因子(图谱网络-动量溢出)}

\subsection{因子定义}

对于任意两只股票 \(i\) 和 \(j\)，计算它们在特征空间中的欧几里得距离：
\begin{equation}
    D_{i,j} = \sqrt{(\mathrm{Prc}_i - \mathrm{Prc}_j)^2 + (\mathrm{SIZE}_i - \mathrm{SIZE}_j)^2 + (\mathrm{BM}_i - \mathrm{BM}_j)^2 + (\mathrm{OP}_i - \mathrm{OP}_j)^2 + (\mathrm{INV}_i - \mathrm{INV}_j)^2}
\end{equation}

对每只股票，将其距离最近的50只股票（最相似的50只股票）的过去一个月的超额回报进行加权平均即为公司特征相似度动量因子：
\begin{equation}
    \mathrm{SIM} = \frac{\sum_{k \in N(i)} w_k \cdot R_{\mathrm{excess},k}}{\sum_{k \in N(i)} w_k}
\end{equation}

\subsection{变量定义}
\begin{itemize}
    \item ① 计算距离的特征分别为：月末收盘价（\(\mathrm{Prc}\)）、对数市值（\(\mathrm{SIZE}\)）、账面市值比（\(\mathrm{BM}\)）、经营利润/净资产（\(\mathrm{OP}\)）和总资产同比增速（\(\mathrm{INV}\)）；
    \item ② \(w_k\) 为股票 \(k\) 的市值权重；
    \item ③ \(R_{\mathrm{excess},k}\) 为股票 \(k\) 在过去一个月的超额收益率。
\end{itemize}

\subsection{说明}
具有相似基本面特征或是相似因子暴露的公司之间有着相似的预期收益，可以利用相似公司的历史收益来预测自身的未来收益；与其相似的那些股票的收益越高，该股票未来收益也越高。

\subsection{参考文献}
\begin{enumerate}
    \item He, W., Wang, Y., \& Yu, J. (2021). \textit{Similar Stocks}.
\end{enumerate}
\end{document}